\documentclass[10pt,a4paper]{amsart}
\usepackage{amsmath,amssymb,amsthm,graphicx}
% For commutative diagrams we use Xy-pic
\usepackage[all]{xy}

% We use 2cell for 2-commutative diagrams.
\xyoption{2cell}
% Theorem environments.
%
\theoremstyle{plain}
\newtheorem{theorem}[subsection]{Theorem}
\newtheorem{proposition}[subsection]{Proposition}
\newtheorem{lemma}[subsection]{Lemma}

\theoremstyle{definition}
\newtheorem{definition}[subsection]{Definition}
\newtheorem{example}[subsection]{Example}
\newtheorem{exercise}[subsection]{Exercise}
\newtheorem{situation}[subsection]{Situation}

\theoremstyle{remark}
\newtheorem{remark}[subsection]{Remark}
\newtheorem{remarks}[subsection]{Remarks}

\numberwithin{equation}{subsection}

% Macros
%
\def\lim{\mathop{\mathrm{lim}}\nolimits}
\def\colim{\mathop{\mathrm{colim}}\nolimits}
\def\Spec{\mathop{\mathrm{Spec}}}
\def\Hom{\mathop{\mathrm{Hom}}\nolimits}
\def\Ext{\mathop{\mathrm{Ext}}\nolimits}
\def\SheafHom{\mathop{\mathcal{H}\!\mathit{om}}\nolimits}
\def\SheafExt{\mathop{\mathcal{E}\!\mathit{xt}}\nolimits}
\def\Sch{\mathit{Sch}}
\def\Mor{\mathop{\mathrm{Mor}}\nolimits}
\def\Ob{\mathop{\mathrm{Ob}}\nolimits}
\def\Sh{\mathop{\mathit{Sh}}\nolimits}
\def\NL{\mathop{N\!L}\nolimits}
\def\CH{\mathop{\mathrm{CH}}\nolimits}
\def\proetale{{pro\text{-}\acute{e}tale}}
\def\etale{{\acute{e}tale}}
\def\QCoh{\mathit{QCoh}}
\def\Ker{\mathop{\mathrm{Ker}}}
\def\Im{\mathop{\mathrm{Im}}}
\def\Coker{\mathop{\mathrm{Coker}}}
\def\Coim{\mathop{\mathrm{Coim}}}

% Boxtimes
%
\DeclareMathSymbol{\boxtimes}{\mathbin}{AMSa}{"02}

%
% Macros for moduli stacks/spaces
%
\def\QCohstack{\mathcal{QC}\!\mathit{oh}}
\def\Cohstack{\mathcal{C}\!\mathit{oh}}
\def\Spacesstack{\mathcal{S}\!\mathit{paces}}
\def\Quotfunctor{\mathrm{Quot}}
\def\Hilbfunctor{\mathrm{Hilb}}
\def\Curvesstack{\mathcal{C}\!\mathit{urves}}
\def\Polarizedstack{\mathcal{P}\!\mathit{olarized}}
\def\Complexesstack{\mathcal{C}\!\mathit{omplexes}}
% \Pic is the operator that assigns to X its picard group, usage \Pic(X)
% \Picardstack_{X/B} denotes the Picard stack of X over B
% \Picardfunctor_{X/B} denotes the Picard functor of X over B
\def\Pic{\mathop{\mathrm{Pic}}\nolimits}
\def\Picardstack{\mathcal{P}\!\mathit{ic}}
\def\Picardfunctor{\mathrm{Pic}}
\def\Deformationcategory{\mathcal{D}\!\mathit{ef}}
\newenvironment{reference}{}{}
\newenvironment{slogan}{}{}
\newenvironment{history}{}{}
\usepackage{cleveref}
\crefname{originallemma}{original lemma}{original lemmas}
\Crefname{originallemma}{Original lemma}{Original lemmas}
\makeatletter
\newlabel{sites-cohomology-lemma-O-homology-qis}{{08RX}{source}{}{sourcepointer-1}{}}
\newlabel{sites-cohomology-lemma-O-homology-qis@cref}{{[originallemma][0][]08RX}{[1][0][]source}{}{}{}}
\makeatother
\crefname{originalsection}{original section}{original sections}
\Crefname{originalsection}{Original section}{Original sections}
\makeatletter
\newlabel{remark-resolution}{{08QI}{source}{}{sourcepointer-2}{}}
\newlabel{remark-resolution@cref}{{[originalsection][0][]08QI}{[1][0][]source}{}{}{}}
\makeatother
\crefname{originalsection}{original section}{original sections}
\Crefname{originalsection}{Original section}{Original sections}
\makeatletter
\newlabel{sites-cohomology-example-category-to-point}{{08PF}{source}{}{sourcepointer-3}{}}
\newlabel{sites-cohomology-example-category-to-point@cref}{{[originalsection][0][]08PF}{[1][0][]source}{}{}{}}
\makeatother
\crefname{originallemma}{original lemma}{original lemmas}
\Crefname{originallemma}{Original lemma}{Original lemmas}
\makeatletter
\newlabel{sites-cohomology-lemma-compute-by-cosimplicial-resolution}{{08Q9}{source}{}{sourcepointer-4}{}}
\newlabel{sites-cohomology-lemma-compute-by-cosimplicial-resolution@cref}{{[originallemma][0][]08Q9}{[1][0][]source}{}{}{}}
\makeatother
\crefname{originallemma}{original lemma}{original lemmas}
\Crefname{originallemma}{Original lemma}{Original lemmas}
\makeatletter
\newlabel{sites-cohomology-lemma-eilenberg-zilber}{{08QC}{source}{}{sourcepointer-5}{}}
\newlabel{sites-cohomology-lemma-eilenberg-zilber@cref}{{[originallemma][0][]08QC}{[1][0][]source}{}{}{}}
\makeatother
\crefname{originallemma}{original lemma}{original lemmas}
\Crefname{originallemma}{Original lemma}{Original lemmas}
\makeatletter
\newlabel{algebra-lemma-trivial-differential-surjective}{{00RP}{source}{}{sourcepointer-6}{}}
\newlabel{algebra-lemma-trivial-differential-surjective@cref}{{[originallemma][0][]00RP}{[1][0][]source}{}{}{}}
\makeatother
\crefname{originallemma}{original lemma}{original lemmas}
\Crefname{originallemma}{Original lemma}{Original lemmas}
\makeatletter
\newlabel{sites-cohomology-lemma-change-of-rings}{{08Q8}{source}{}{sourcepointer-7}{}}
\newlabel{sites-cohomology-lemma-change-of-rings@cref}{{[originallemma][0][]08Q8}{[1][0][]source}{}{}{}}
\makeatother
\crefname{originalsection}{original section}{original sections}
\Crefname{originalsection}{Original section}{Original sections}
\makeatletter
\newlabel{homology-section-filtered-complex}{{012K}{source}{}{sourcepointer-8}{}}
\newlabel{homology-section-filtered-complex@cref}{{[originalsection][0][]012K}{[1][0][]source}{}{}{}}
\makeatother
\crefname{originallemma}{original lemma}{original lemmas}
\Crefname{originallemma}{Original lemma}{Original lemmas}
\makeatletter
\newlabel{derived-lemma-trick-vanishing-composition}{{08Q2}{source}{}{sourcepointer-9}{}}
\newlabel{derived-lemma-trick-vanishing-composition@cref}{{[originallemma][0][]08Q2}{[1][0][]source}{}{}{}}
\makeatother
\crefname{originalsection}{original section}{original sections}
\Crefname{originalsection}{Original section}{Original sections}
\makeatletter
\newlabel{more-algebra-example-tor-change-rings}{{068F}{source}{}{sourcepointer-10}{}}
\newlabel{more-algebra-example-tor-change-rings@cref}{{[originalsection][0][]068F}{[1][0][]source}{}{}{}}
\makeatother
\crefname{originalsection}{original section}{original sections}
\Crefname{originalsection}{Original section}{Original sections}
\makeatletter
\newlabel{simplicial-section-standard}{{08N8}{source}{}{sourcepointer-11}{}}
\newlabel{simplicial-section-standard@cref}{{[originalsection][0][]08N8}{[1][0][]source}{}{}{}}
\makeatother
\crefname{originalsection}{original section}{original sections}
\Crefname{originalsection}{Original section}{Original sections}
\makeatletter
\newlabel{simplicial-example-polynomial-algebra-maps}{{09CB}{source}{}{sourcepointer-12}{}}
\newlabel{simplicial-example-polynomial-algebra-maps@cref}{{[originalsection][0][]09CB}{[1][0][]source}{}{}{}}
\makeatother
\crefname{originalsection}{original section}{original sections}
\Crefname{originalsection}{Original section}{Original sections}
\makeatletter
\newlabel{simplicial-section-complexes}{{0194}{source}{}{sourcepointer-13}{}}
\newlabel{simplicial-section-complexes@cref}{{[originalsection][0][]0194}{[1][0][]source}{}{}{}}
\makeatother
\crefname{originallemma}{original lemma}{original lemmas}
\Crefname{originallemma}{Original lemma}{Original lemmas}
\makeatletter
\newlabel{algebra-lemma-colimit-differentials}{{031G}{source}{}{sourcepointer-14}{}}
\newlabel{algebra-lemma-colimit-differentials@cref}{{[originallemma][0][]031G}{[1][0][]source}{}{}{}}
\makeatother
\crefname{originalsection}{original section}{original sections}
\Crefname{originalsection}{Original section}{Original sections}
\makeatletter
\newlabel{sites-example-indiscrete}{{07GE}{source}{}{sourcepointer-15}{}}
\newlabel{sites-example-indiscrete@cref}{{[originalsection][0][]07GE}{[1][0][]source}{}{}{}}
\makeatother
\crefname{originallemma}{original lemma}{original lemmas}
\Crefname{originallemma}{Original lemma}{Original lemmas}
\makeatletter
\newlabel{simplicial-lemma-trivial-kan-base-change}{{08NN}{source}{}{sourcepointer-16}{}}
\newlabel{simplicial-lemma-trivial-kan-base-change@cref}{{[originallemma][0][]08NN}{[1][0][]source}{}{}{}}
\makeatother
\crefname{originallemma}{original lemma}{original lemmas}
\Crefname{originallemma}{Original lemma}{Original lemmas}
\makeatletter
\newlabel{simplicial-lemma-product-trivial-kan}{{08NR}{source}{}{sourcepointer-17}{}}
\newlabel{simplicial-lemma-product-trivial-kan@cref}{{[originallemma][0][]08NR}{[1][0][]source}{}{}{}}
\makeatother
\crefname{originallemma}{original lemma}{original lemmas}
\Crefname{originallemma}{Original lemma}{Original lemmas}
\makeatletter
\newlabel{simplicial-lemma-trivial-kan-homotopy}{{08NS}{source}{}{sourcepointer-18}{}}
\newlabel{simplicial-lemma-trivial-kan-homotopy@cref}{{[originallemma][0][]08NS}{[1][0][]source}{}{}{}}
\makeatother
\crefname{originallemma}{original lemma}{original lemmas}
\Crefname{originallemma}{Original lemma}{Original lemmas}
\makeatletter
\newlabel{simplicial-lemma-standard-simplicial-homotopy}{{08ND}{source}{}{sourcepointer-19}{}}
\newlabel{simplicial-lemma-standard-simplicial-homotopy@cref}{{[originallemma][0][]08ND}{[1][0][]source}{}{}{}}
\makeatother
\crefname{originallemma}{original lemma}{original lemmas}
\Crefname{originallemma}{Original lemma}{Original lemmas}
\makeatletter
\newlabel{simplicial-lemma-homotopy-equivalence}{{08P2}{source}{}{sourcepointer-20}{}}
\newlabel{simplicial-lemma-homotopy-equivalence@cref}{{[originallemma][0][]08P2}{[1][0][]source}{}{}{}}
\makeatother
\crefname{originallemma}{original lemma}{original lemmas}
\Crefname{originallemma}{Original lemma}{Original lemmas}
\makeatletter
\newlabel{simplicial-lemma-qis-simplicial-abelian-groups}{{08P1}{source}{}{sourcepointer-21}{}}
\newlabel{simplicial-lemma-qis-simplicial-abelian-groups@cref}{{[originallemma][0][]08P1}{[1][0][]source}{}{}{}}
\makeatother
\crefname{originallemma}{original lemma}{original lemmas}
\Crefname{originallemma}{Original lemma}{Original lemmas}
\makeatletter
\newlabel{algebra-lemma-differential-surjective}{{00RR}{source}{}{sourcepointer-22}{}}
\newlabel{algebra-lemma-differential-surjective@cref}{{[originallemma][0][]00RR}{[1][0][]source}{}{}{}}
\makeatother
\crefname{originallemma}{original lemma}{original lemmas}
\Crefname{originallemma}{Original lemma}{Original lemmas}
\makeatletter
\newlabel{algebra-lemma-differential-seq}{{00RU}{source}{}{sourcepointer-23}{}}
\newlabel{algebra-lemma-differential-seq@cref}{{[originallemma][0][]00RU}{[1][0][]source}{}{}{}}
\makeatother
\title{Quillen spectral sequence from symmetric powers of cotangent data to Tor}
\author{The Stacks Project Authors}
\date{stacks-project:cotangent.tex:local-snapshot-20261011; GNU FDL 1.2 or later}
\begin{document}
\maketitle
\noindent\textit{Editorial note.} The09CG task identifies a three-term truncation of the cotangent complex. This task constructs an all-degree spectral sequence and proves its convergence; neither statement is a relabelled proof of the other.\par\medskip
\noindent\textit{Allowed standard prerequisites (editorial).}\begin{itemize}
\item Explicit prerequisite boundary. Simplicial polynomial resolutions, derived lower shriek, filtered-complex spectral sequences, change-of-ringsTor and truncation composition. The local comparison, symmetric-power vanishing and Tor-map-zero proofs are retained.
\end{itemize}\medskip
\ifdefined\hypertarget\hypertarget{sourcepointer-1}{}\fi
\noindent\textit{Source reference.} Original lemma 08RX: focused-seeds/AGP/sources/stacks-snapshot-20261011/tags.txt, source lines 11110--11110. Stable Stacks tag for an expressly allowed foundational or notational input; the counted construction remains the two concrete core obligations. The retained local core is not replaced by this pointer. See standard\_prerequisites.\par\medskip
\ifdefined\hypertarget\hypertarget{sourcepointer-2}{}\fi
\noindent\textit{Source reference.} Original section 08QI: focused-seeds/AGP/sources/stacks-snapshot-20261011/tags.txt, source lines 11061--11061. Stable Stacks tag for an expressly allowed foundational or notational input; the counted construction remains the two concrete core obligations. The retained local core is not replaced by this pointer. See standard\_prerequisites.\par\medskip
\ifdefined\hypertarget\hypertarget{sourcepointer-3}{}\fi
\noindent\textit{Source reference.} Original section 08PF: focused-seeds/AGP/sources/stacks-snapshot-20261011/tags.txt, source lines 11023--11023. Stable Stacks tag for an expressly allowed foundational or notational input; the counted construction remains the two concrete core obligations. The retained local core is not replaced by this pointer. See standard\_prerequisites.\par\medskip
\ifdefined\hypertarget\hypertarget{sourcepointer-4}{}\fi
\noindent\textit{Source reference.} Original lemma 08Q9: focused-seeds/AGP/sources/stacks-snapshot-20261011/tags.txt, source lines 11052--11052. Stable Stacks tag for an expressly allowed foundational or notational input; the counted construction remains the two concrete core obligations. The retained local core is not replaced by this pointer. See standard\_prerequisites.\par\medskip
\ifdefined\hypertarget\hypertarget{sourcepointer-5}{}\fi
\noindent\textit{Source reference.} Original lemma 08QC: focused-seeds/AGP/sources/stacks-snapshot-20261011/tags.txt, source lines 11055--11055. Stable Stacks tag for an expressly allowed foundational or notational input; the counted construction remains the two concrete core obligations. The retained local core is not replaced by this pointer. See standard\_prerequisites.\par\medskip
\ifdefined\hypertarget\hypertarget{sourcepointer-6}{}\fi
\noindent\textit{Source reference.} Original lemma 00RP: focused-seeds/AGP/sources/stacks-snapshot-20261011/tags.txt, source lines 1011--1011. Stable Stacks tag for an expressly allowed foundational or notational input; the counted construction remains the two concrete core obligations. The retained local core is not replaced by this pointer. See standard\_prerequisites.\par\medskip
\ifdefined\hypertarget\hypertarget{sourcepointer-7}{}\fi
\noindent\textit{Source reference.} Original lemma 08Q8: focused-seeds/AGP/sources/stacks-snapshot-20261011/tags.txt, source lines 11051--11051. Stable Stacks tag for an expressly allowed foundational or notational input; the counted construction remains the two concrete core obligations. The retained local core is not replaced by this pointer. See standard\_prerequisites.\par\medskip
\ifdefined\hypertarget\hypertarget{sourcepointer-8}{}\fi
\noindent\textit{Source reference.} Original section 012K: focused-seeds/AGP/sources/stacks-snapshot-20261011/tags.txt, source lines 1402--1402. Stable Stacks tag for an expressly allowed foundational or notational input; the counted construction remains the two concrete core obligations. The retained local core is not replaced by this pointer. See standard\_prerequisites.\par\medskip
\ifdefined\hypertarget\hypertarget{sourcepointer-9}{}\fi
\noindent\textit{Source reference.} Original lemma 08Q2: focused-seeds/AGP/sources/stacks-snapshot-20261011/tags.txt, source lines 11045--11045. Stable Stacks tag for an expressly allowed foundational or notational input; the counted construction remains the two concrete core obligations. The retained local core is not replaced by this pointer. See standard\_prerequisites.\par\medskip
\ifdefined\hypertarget\hypertarget{sourcepointer-10}{}\fi
\noindent\textit{Source reference.} Original section 068F: focused-seeds/AGP/sources/stacks-snapshot-20261011/tags.txt, source lines 8004--8004. Stable Stacks tag for an expressly allowed foundational or notational input; the counted construction remains the two concrete core obligations. The retained local core is not replaced by this pointer. See standard\_prerequisites.\par\medskip
\ifdefined\hypertarget\hypertarget{sourcepointer-11}{}\fi
\noindent\textit{Source reference.} Original section 08N8: focused-seeds/AGP/sources/stacks-snapshot-20261011/tags.txt, source lines 10981--10981. Stable Stacks tag for an expressly allowed foundational or notational input; the counted construction remains the two concrete core obligations. The retained local core is not replaced by this pointer. See standard\_prerequisites.\par\medskip
\ifdefined\hypertarget\hypertarget{sourcepointer-12}{}\fi
\noindent\textit{Source reference.} Original section 09CB: focused-seeds/AGP/sources/stacks-snapshot-20261011/tags.txt, source lines 11824--11824. Stable Stacks tag for an expressly allowed foundational or notational input; the counted construction remains the two concrete core obligations. The retained local core is not replaced by this pointer. See standard\_prerequisites.\par\medskip
\ifdefined\hypertarget\hypertarget{sourcepointer-13}{}\fi
\noindent\textit{Source reference.} Original section 0194: focused-seeds/AGP/sources/stacks-snapshot-20261011/tags.txt, source lines 1639--1639. Stable Stacks tag for an expressly allowed foundational or notational input; the counted construction remains the two concrete core obligations. The retained local core is not replaced by this pointer. See standard\_prerequisites.\par\medskip
\ifdefined\hypertarget\hypertarget{sourcepointer-14}{}\fi
\noindent\textit{Source reference.} Original lemma 031G: focused-seeds/AGP/sources/stacks-snapshot-20261011/tags.txt, source lines 3983--3983. Stable Stacks tag for an expressly allowed foundational or notational input; the counted construction remains the two concrete core obligations. The retained local core is not replaced by this pointer. See standard\_prerequisites.\par\medskip
\ifdefined\hypertarget\hypertarget{sourcepointer-15}{}\fi
\noindent\textit{Source reference.} Original section 07GE: focused-seeds/AGP/sources/stacks-snapshot-20261011/tags.txt, source lines 9514--9514. Stable Stacks tag for an expressly allowed foundational or notational input; the counted construction remains the two concrete core obligations. The retained local core is not replaced by this pointer. See standard\_prerequisites.\par\medskip
\ifdefined\hypertarget\hypertarget{sourcepointer-16}{}\fi
\noindent\textit{Source reference.} Original lemma 08NN: focused-seeds/AGP/sources/stacks-snapshot-20261011/tags.txt, source lines 10996--10996. Stable Stacks tag for an expressly allowed foundational or notational input; the counted construction remains the two concrete core obligations. The retained local core is not replaced by this pointer. See standard\_prerequisites.\par\medskip
\ifdefined\hypertarget\hypertarget{sourcepointer-17}{}\fi
\noindent\textit{Source reference.} Original lemma 08NR: focused-seeds/AGP/sources/stacks-snapshot-20261011/tags.txt, source lines 10999--10999. Stable Stacks tag for an expressly allowed foundational or notational input; the counted construction remains the two concrete core obligations. The retained local core is not replaced by this pointer. See standard\_prerequisites.\par\medskip
\ifdefined\hypertarget\hypertarget{sourcepointer-18}{}\fi
\noindent\textit{Source reference.} Original lemma 08NS: focused-seeds/AGP/sources/stacks-snapshot-20261011/tags.txt, source lines 11000--11000. Stable Stacks tag for an expressly allowed foundational or notational input; the counted construction remains the two concrete core obligations. The retained local core is not replaced by this pointer. See standard\_prerequisites.\par\medskip
\ifdefined\hypertarget\hypertarget{sourcepointer-19}{}\fi
\noindent\textit{Source reference.} Original lemma 08ND: focused-seeds/AGP/sources/stacks-snapshot-20261011/tags.txt, source lines 10986--10986. Stable Stacks tag for an expressly allowed foundational or notational input; the counted construction remains the two concrete core obligations. The retained local core is not replaced by this pointer. See standard\_prerequisites.\par\medskip
\ifdefined\hypertarget\hypertarget{sourcepointer-20}{}\fi
\noindent\textit{Source reference.} Original lemma 08P2: focused-seeds/AGP/sources/stacks-snapshot-20261011/tags.txt, source lines 11010--11010. Stable Stacks tag for an expressly allowed foundational or notational input; the counted construction remains the two concrete core obligations. The retained local core is not replaced by this pointer. See standard\_prerequisites.\par\medskip
\ifdefined\hypertarget\hypertarget{sourcepointer-21}{}\fi
\noindent\textit{Source reference.} Original lemma 08P1: focused-seeds/AGP/sources/stacks-snapshot-20261011/tags.txt, source lines 11009--11009. Stable Stacks tag for an expressly allowed foundational or notational input; the counted construction remains the two concrete core obligations. The retained local core is not replaced by this pointer. See standard\_prerequisites.\par\medskip
\ifdefined\hypertarget\hypertarget{sourcepointer-22}{}\fi
\noindent\textit{Source reference.} Original lemma 00RR: focused-seeds/AGP/sources/stacks-snapshot-20261011/tags.txt, source lines 1013--1013. Stable Stacks tag for an expressly allowed foundational or notational input; the counted construction remains the two concrete core obligations. The retained local core is not replaced by this pointer. See standard\_prerequisites.\par\medskip
\ifdefined\hypertarget\hypertarget{sourcepointer-23}{}\fi
\noindent\textit{Source reference.} Original lemma 00RU: focused-seeds/AGP/sources/stacks-snapshot-20261011/tags.txt, source lines 1016--1016. Stable Stacks tag for an expressly allowed foundational or notational input; the counted construction remains the two concrete core obligations. The retained local core is not replaced by this pointer. See standard\_prerequisites.\par\medskip
\begin{theorem}[Quillen spectral sequence]
\label{theorem-quillen-spectral-sequence}
Let $A \to B$ be a surjective ring map. Consider the sheaf
$\Omega = \Omega_{\mathcal{O}/A} \otimes_\mathcal{O} \underline{B}$ of
$\underline{B}$-modules on $\mathcal{C}_{B/A}$, see
Section \ref{section-compute-L-pi-shriek}.
Then there is a spectral sequence with $E_1$-page
$$
E_1^{p, q} =
H_{- p - q}(\mathcal{C}_{B/A}, \text{Sym}^p_{\underline{B}}(\Omega))
\Rightarrow \text{Tor}^A_{- p - q}(B, B)
$$
with $d_r$ of bidegree $(r, -r + 1)$.
Moreover, $H_i(\mathcal{C}_{B/A}, \text{Sym}^k_{\underline{B}}(\Omega)) = 0$
for $i < k$.
\end{theorem}

\section{The cotangent complex of a ring map}
\label{section-cotangent-ring-map}

\noindent
Let $A$ be a ring. Let $\textit{Alg}_A$ be the category of $A$-algebras.
Consider the pair of adjoint functors $(U, V)$ where
$V : \textit{Alg}_A \to \textit{Sets}$ is the forgetful functor and
$U : \textit{Sets} \to \textit{Alg}_A$ assigns to a set $E$ the polynomial
algebra $A[E]$ on $E$ over $A$. Let $X_\bullet$ be the simplicial object of
$\text{Fun}(\textit{Alg}_A, \textit{Alg}_A)$ constructed in
Simplicial, Section \ref{simplicial-section-standard}.

\medskip\noindent
Consider an $A$-algebra $B$. Denote $P_\bullet = X_\bullet(B)$ the resulting
simplicial $A$-algebra. Recall that $P_0 = A[B]$, $P_1 = A[A[B]]$, and so on.
In particular each term $P_n$ is a polynomial $A$-algebra.
Recall also that there is an augmentation
$$
\epsilon : P_\bullet \longrightarrow B
$$
where we view $B$ as a constant simplicial $A$-algebra.

\begin{definition}
\label{definition-standard-resolution}
Let $A \to B$ be a ring map. The {\it standard resolution of $B$ over $A$}
is the augmentation $\epsilon : P_\bullet \to B$ with terms
$$
P_0 = A[B],\quad P_1 = A[A[B]],\quad \ldots
$$
and maps as constructed in
Simplicial, Example \ref{simplicial-example-polynomial-algebra-maps}.
\end{definition}

\noindent
It will turn out that we can use the standard resolution
to compute left derived functors in certain settings.

\begin{definition}
\label{definition-cotangent-complex-ring-map}
The {\it cotangent complex} $L_{B/A}$ of a ring map $A \to B$
is the complex of $B$-modules associated to the simplicial $B$-module
$$
\Omega_{P_\bullet/A} \otimes_{P_\bullet, \epsilon} B
$$
where $\epsilon : P_\bullet \to B$ is the standard resolution
of $B$ over $A$.
\end{definition}

\noindent
In Simplicial, Section \ref{simplicial-section-complexes} we associate a
chain complex to a simplicial module, but here we work with cochain complexes.
Thus the term $L_{B/A}^{-n}$ in degree $-n$ is the $B$-module
$\Omega_{P_n/A} \otimes_{P_n, \epsilon_n} B$ and $L_{B/A}^m = 0$
for $m > 0$.

\begin{remark}
\label{remark-variant-cotangent-complex}
Let $A \to B$ be a ring map. Let $\mathcal{A}$ be the category of
arrows $\psi : C \to B$ of $A$-algebras and let $\mathcal{S}$ be
the category of maps $E \to B$ where $E$ is a set. There are adjoint
functors $V : \mathcal{A} \to \mathcal{S}$ (the forgetful functor)
and $U : \mathcal{S} \to \mathcal{A}$ which sends $E \to B$ to
$A[E] \to B$. Let $X_\bullet$ be the simplicial object of
$\text{Fun}(\mathcal{A}, \mathcal{A})$ constructed in
Simplicial, Section \ref{simplicial-section-standard}.
The diagram
$$
\xymatrix{
\mathcal{A} \ar[d] \ar[r] & \mathcal{S} \ar@<1ex>[l] \ar[d] \\
\textit{Alg}_A \ar[r] & \textit{Sets} \ar@<1ex>[l]
}
$$
commutes. It follows that $X_\bullet(\text{id}_B : B \to B)$
is equal to the standard resolution of $B$ over $A$.
\end{remark}

\begin{lemma}
\label{lemma-colimit-cotangent-complex}
Let $A_i \to B_i$ be a system of ring maps over a directed index
set $I$. Then $\colim L_{B_i/A_i} = L_{\colim B_i/\colim A_i}$.
\end{lemma}

\begin{proof}
This is true because the forgetful functor
$V : A\textit{-Alg} \to \textit{Sets}$ and its adjoint
$U : \textit{Sets} \to A\textit{-Alg}$ commute with filtered colimits.
Moreover, the functor $B/A \mapsto \Omega_{B/A}$ does as well
(Algebra, Lemma \ref{algebra-lemma-colimit-differentials}).
\end{proof}





\section{Simplicial resolutions and derived lower shriek}
\label{section-compute-L-pi-shriek}

\noindent
Let $A \to B$ be a ring map. Consider the category whose objects are
$A$-algebra maps $\alpha : P \to B$ where $P$ is a polynomial algebra over $A$
(in some set\footnote{It suffices to consider sets of cardinality
at most the cardinality of $B$.} of variables) and whose
morphisms $s : (\alpha : P \to B) \to (\alpha' : P' \to B)$ are
$A$-algebra homomorphisms $s : P \to P'$ with $\alpha' \circ s = \alpha$.
Let $\mathcal{C} = \mathcal{C}_{B/A}$ denote the {\bf opposite}
of this category. The reason for
taking the opposite is that we want to think of objects
$(P, \alpha)$ as corresponding to the diagram of affine schemes
$$
\xymatrix{
\Spec(B) \ar[d] \ar[r] & \Spec(P) \ar[ld] \\
\Spec(A)
}
$$
We endow $\mathcal{C}$ with the chaotic topology
(Sites, Example \ref{sites-example-indiscrete}), i.e., we endow
$\mathcal{C}$ with the structure of a site where coverings are given by
identities so that all presheaves are sheaves.
Moreover, we endow $\mathcal{C}$ with two sheaves of rings. The first
is the sheaf $\mathcal{O}$ which sends to object $(P, \alpha)$ to $P$.
Then second is the constant sheaf $B$, which we will denote
$\underline{B}$. We obtain the following diagram of morphisms of
ringed topoi
\begin{equation}
\label{equation-pi}
\vcenter{
\xymatrix{
(\Sh(\mathcal{C}), \underline{B}) \ar[r]_i \ar[d]_\pi &
(\Sh(\mathcal{C}), \mathcal{O}) \\
(\Sh(*), B)
}
}
\end{equation}
The morphism $i$ is the identity on underlying topoi and
$i^\sharp : \mathcal{O} \to \underline{B}$ is the obvious map.
The map $\pi$ is as in Cohomology on Sites, Example
\ref{sites-cohomology-example-category-to-point}.
An important role will be played in the following
by the derived functors
$
Li^* : D(\mathcal{O}) \longrightarrow D(\underline{B})
$
left adjoint to $Ri_* = i_* : D(\underline{B}) \to D(\mathcal{O})$ and
$
L\pi_! : D(\underline{B}) \longrightarrow D(B)
$
left adjoint to $\pi^* = \pi^{-1} : D(B) \to D(\underline{B})$.

\begin{lemma}
\label{lemma-identify-pi-shriek}
With notation as above let $P_\bullet$ be a simplicial $A$-algebra
endowed with an augmentation $\epsilon : P_\bullet \to B$.
Assume each $P_n$ is a polynomial algebra over $A$ and $\epsilon$
is a trivial Kan fibration on underlying simplicial sets. Then
$$
L\pi_!(\mathcal{F}) = \mathcal{F}(P_\bullet, \epsilon)
$$
in $D(\textit{Ab})$, resp.\ $D(B)$ functorially in $\mathcal{F}$ in
$\textit{Ab}(\mathcal{C})$, resp.\ $\textit{Mod}(\underline{B})$.
\end{lemma}

\begin{proof}
We will use the criterion of Cohomology on Sites, Lemma
\ref{sites-cohomology-lemma-compute-by-cosimplicial-resolution} to prove this.
Given an object $U = (Q, \beta)$ of $\mathcal{C}$ we have to show that
$$
S_\bullet = \Mor_\mathcal{C}((P_\bullet, \epsilon), (Q, \beta))
$$
is homotopy equivalent to a singleton.
Write $Q = A[E]$ for some set $E$ (this is possible by our choice of
the category $\mathcal{C}$). We see that
$$
S_\bullet = \Mor_{\textit{Sets}}((E, \beta|_E), (P_\bullet, \epsilon))
$$
Let $*$ be the constant simplicial set on a singleton. For $b \in B$
let $F_{b, \bullet}$ be the simplicial set defined by the cartesian
diagram
$$
\xymatrix{
F_{b, \bullet} \ar[r] \ar[d] & P_\bullet \ar[d]_\epsilon \\
{*} \ar[r]^b & B
}
$$
With this notation $S_\bullet = \prod_{e \in E} F_{\beta(e), \bullet}$.
Since we assumed $\epsilon$ is a trivial Kan fibration we see that
$F_{b, \bullet} \to *$ is a trivial Kan fibration
(Simplicial, Lemma \ref{simplicial-lemma-trivial-kan-base-change}).
Thus $S_\bullet \to *$ is a trivial Kan fibration
(Simplicial, Lemma \ref{simplicial-lemma-product-trivial-kan}).
Therefore $S_\bullet$ is homotopy equivalent to $*$
(Simplicial, Lemma \ref{simplicial-lemma-trivial-kan-homotopy}).
\end{proof}

\noindent
In particular, we can use the standard resolution of $B$ over $A$
to compute derived lower shriek.

\begin{lemma}
\label{lemma-pi-shriek-standard}
Let $A \to B$ be a ring map. Let $\epsilon : P_\bullet \to B$
be the standard resolution of $B$ over $A$. Let $\pi$ be as in
(\ref{equation-pi}). Then
$$
L\pi_!(\mathcal{F}) = \mathcal{F}(P_\bullet, \epsilon)
$$
in $D(\textit{Ab})$, resp.\ $D(B)$ functorially in $\mathcal{F}$ in
$\textit{Ab}(\mathcal{C})$, resp.\ $\textit{Mod}(\underline{B})$.
\end{lemma}

\begin{proof}[First proof]
We will apply Lemma \ref{lemma-identify-pi-shriek}.
Since the terms $P_n$ are polynomial algebras we see the first
assumption of that lemma is satisfied. The second assumption is proved
as follows. By
Simplicial, Lemma \ref{simplicial-lemma-standard-simplicial-homotopy}
the map $\epsilon$ is a homotopy equivalence of underlying
simplicial sets. By
Simplicial, Lemma \ref{simplicial-lemma-homotopy-equivalence}
this implies $\epsilon$ induces a quasi-isomorphism of associated
complexes of abelian groups. By
Simplicial, Lemma \ref{simplicial-lemma-qis-simplicial-abelian-groups}
this implies that $\epsilon$ is a trivial Kan fibration of underlying
simplicial sets.
\end{proof}

\begin{proof}[Second proof]
We will use the criterion of Cohomology on Sites, Lemma
\ref{sites-cohomology-lemma-compute-by-cosimplicial-resolution}.
Let $U = (Q, \beta)$ be an object of $\mathcal{C}$.
We have to show that
$$
S_\bullet = \Mor_\mathcal{C}((P_\bullet, \epsilon), (Q, \beta))
$$
is homotopy equivalent to a singleton. Write $Q = A[E]$ for some set $E$
(this is possible by our choice of the category $\mathcal{C}$). Using the
notation of Remark \ref{remark-variant-cotangent-complex} we see that
$$
S_\bullet = \Mor_\mathcal{S}((E \to B), i(P_\bullet \to B))
$$
By Simplicial, Lemma \ref{simplicial-lemma-standard-simplicial-homotopy}
the map $i(P_\bullet \to B) \to i(B \to B)$ is a homotopy equivalence
in $\mathcal{S}$. Hence $S_\bullet$ is homotopy equivalent to
$$
\Mor_\mathcal{S}((E \to B), (B \to B)) = \{*\}
$$
as desired.
\end{proof}

\begin{lemma}
\label{lemma-compute-cotangent-complex}
Let $A \to B$ be a ring map.  Let $\pi$ and $i$ be as in (\ref{equation-pi}).
There is a canonical isomorphism
$$
L_{B/A} = L\pi_!(Li^*\Omega_{\mathcal{O}/A}) =
L\pi_!(i^*\Omega_{\mathcal{O}/A}) =
L\pi_!(\Omega_{\mathcal{O}/A} \otimes_\mathcal{O} \underline{B})
$$
in $D(B)$.
\end{lemma}

\begin{proof}
For an object $\alpha : P \to B$ of the category $\mathcal{C}$
the module $\Omega_{P/A}$ is a free $P$-module. Thus
$\Omega_{\mathcal{O}/A}$ is a flat $\mathcal{O}$-module. Hence
$Li^*\Omega_{\mathcal{O}/A} = i^*\Omega_{\mathcal{O}/A}$ is the sheaf
of $\underline{B}$-modules which associates to $\alpha : P \to A$ the
$B$-module $\Omega_{P/A} \otimes_{P, \alpha} B$.
By Lemma \ref{lemma-pi-shriek-standard}
we see that the right hand side is computed by
the value of this sheaf on the standard resolution which is our
definition of the left hand side
(Definition \ref{definition-cotangent-complex-ring-map}).
\end{proof}

\begin{lemma}
\label{lemma-pi-lower-shriek-constant-sheaf}
If $A \to B$ is a ring map, then $L\pi_!(\pi^{-1}M) = M$
with $\pi$ as in (\ref{equation-pi}).
\end{lemma}

\begin{proof}
This follows from Lemma \ref{lemma-identify-pi-shriek} which tells us
$L\pi_!(\pi^{-1}M)$ is computed by $(\pi^{-1}M)(P_\bullet, \epsilon)$
which is the constant simplicial object on $M$.
\end{proof}

\begin{lemma}
\label{lemma-identify-H0}
If $A \to B$ is a ring map, then $H^0(L_{B/A}) = \Omega_{B/A}$.
\end{lemma}

\begin{proof}
We will prove this by a direct calculation.
We will use the identification of Lemma \ref{lemma-compute-cotangent-complex}.
There is clearly a map from $\Omega_{\mathcal{O}/A} \otimes \underline{B}$
to the constant sheaf with value $\Omega_{B/A}$. Thus this map induces
a map
$$
H^0(L_{B/A}) = H^0(L\pi_!(\Omega_{\mathcal{O}/A} \otimes \underline{B}))
= \pi_!(\Omega_{\mathcal{O}/A} \otimes \underline{B}) \to \Omega_{B/A}
$$
By choosing an object $P \to B$ of $\mathcal{C}_{B/A}$ with $P \to B$
surjective we see that this map is surjective (by
Algebra, Lemma \ref{algebra-lemma-differential-surjective}).
To show that it is injective, suppose that $P \to B$ is an object
of $\mathcal{C}_{B/A}$ and that $\xi \in \Omega_{P/A} \otimes_P B$
is an element which maps to zero in $\Omega_{B/A}$.
We first choose factorization $P \to P' \to B$ such that $P' \to B$
is surjective and $P'$ is a polynomial algebra over $A$.
We may replace $P$ by $P'$. If $B = P/I$, then the kernel
$\Omega_{P/A} \otimes_P B \to \Omega_{B/A}$ is the image of
$I/I^2$ (Algebra, Lemma \ref{algebra-lemma-differential-seq}).
Say $\xi$ is the image of $f \in I$.
Then we consider the two maps $a, b : P' = P[x] \to P$, the first of which
maps $x$ to $0$ and the second of which maps $x$ to $f$ (in both
cases $P[x] \to B$ maps $x$ to zero). We see that $\xi$ and $0$
are the image of $\text{d}x \otimes 1$ in $\Omega_{P'/A} \otimes_{P'} B$.
Thus $\xi$ and $0$ have the same image in the colimit (see
Cohomology on Sites, Example \ref{sites-cohomology-example-category-to-point})
$\pi_!(\Omega_{\mathcal{O}/A} \otimes \underline{B})$ as desired.
\end{proof}

\begin{lemma}
\label{lemma-pi-lower-shriek-polynomial-algebra}
If $B$ is a polynomial algebra over the ring $A$, then
with $\pi$ as in (\ref{equation-pi}) we have that
$\pi_!$ is exact and $\pi_!\mathcal{F} = \mathcal{F}(B \to B)$.
\end{lemma}

\begin{proof}
This follows from Lemma \ref{lemma-identify-pi-shriek} which tells us
the constant simplicial algebra on $B$ can be used to compute $L\pi_!$.
\end{proof}

\begin{lemma}
\label{lemma-cotangent-complex-polynomial-algebra}
If $B$ is a polynomial algebra over the ring $A$, then
$L_{B/A}$ is quasi-isomorphic to $\Omega_{B/A}[0]$.
\end{lemma}

\begin{proof}
Immediate from
Lemmas \ref{lemma-compute-cotangent-complex} and
\ref{lemma-pi-lower-shriek-polynomial-algebra}.
\end{proof}






\begin{lemma}
\label{lemma-O-homology-B-homology}
Let $A \to B$ be a ring map. Let $\pi$, $\mathcal{O}$, $\underline{B}$
be as in (\ref{equation-pi}). For any $\mathcal{O}$-module $\mathcal{F}$
we have
$$
L\pi_!(\mathcal{F}) = L\pi_!(Li^*\mathcal{F}) =
L\pi_!(\mathcal{F} \otimes_\mathcal{O}^\mathbf{L} \underline{B})
$$
in $D(\textit{Ab})$.
\end{lemma}

\begin{proof}
It suffices to verify the assumptions of Cohomology on Sites, Lemma
\ref{sites-cohomology-lemma-O-homology-qis}
hold for $\mathcal{O} \to \underline{B}$ on $\mathcal{C}_{B/A}$.
We will use the results of Remark \ref{remark-resolution} without
further mention. Choose a resolution $P_\bullet$ of $B$ over $A$ to get a
suitable cosimplicial object $U_\bullet$ of $\mathcal{C}_{B/A}$.
Since $P_\bullet \to B$ induces a quasi-isomorphism on associated
complexes of abelian groups we see that $L\pi_!\mathcal{O} = B$.
On the other hand $L\pi_!\underline{B}$ is computed by
$\underline{B}(U_\bullet) = B$. This verifies the second assumption of
Cohomology on Sites, Lemma
\ref{sites-cohomology-lemma-O-homology-qis}
and we are done with the proof.
\end{proof}

\begin{lemma}
\label{lemma-apply-O-B-comparison}
Let $A \to B$ be a ring map. Let $\pi$, $\mathcal{O}$, $\underline{B}$
be as in (\ref{equation-pi}). We have
$$
L\pi_!(\mathcal{O}) = L\pi_!(\underline{B}) = B
\quad\text{and}\quad
L_{B/A} = L\pi_!(\Omega_{\mathcal{O}/A} \otimes_\mathcal{O} \underline{B}) =
L\pi_!(\Omega_{\mathcal{O}/A})
$$
in $D(\textit{Ab})$.
\end{lemma}

\begin{proof}
This is just an application of Lemma \ref{lemma-O-homology-B-homology}
(and the first equality on the right is
Lemma \ref{lemma-compute-cotangent-complex}).
\end{proof}

\begin{lemma}
\label{lemma-vanishing-symmetric-powers}
Notation and assumptions as in
Cohomology on Sites, Example \ref{sites-cohomology-example-category-to-point}.
Assume $\mathcal{C}$ has a cosimplicial object as in
Cohomology on Sites, Lemma
\ref{sites-cohomology-lemma-compute-by-cosimplicial-resolution}.
Let $\mathcal{F}$ be a flat $\underline{B}$-module such that
$H_0(\mathcal{C}, \mathcal{F}) = 0$.
Then $H_l(\mathcal{C}, \text{Sym}_{\underline{B}}^k(\mathcal{F})) = 0$
for $l < k$.
\end{lemma}

\begin{proof}
We drop the subscript ${}_{\underline{B}}$ from tensor products, wedge powers,
and symmetric powers. We will prove the lemma by induction on $k$.
The cases $k = 0, 1$ follow from the assumptions. If $k > 1$ consider
the exact complex
$$
\ldots \to
\wedge^2\mathcal{F} \otimes \text{Sym}^{k - 2}\mathcal{F} \to
\mathcal{F} \otimes \text{Sym}^{k - 1}\mathcal{F} \to
\text{Sym}^k\mathcal{F} \to 0
$$
with differentials as in the Koszul complex. If we think of this as a
resolution of $\text{Sym}^k\mathcal{F}$, then this gives a first quadrant
spectral sequence
$$
E_1^{p, q} =
H_p(\mathcal{C},
\wedge^{q + 1}\mathcal{F} \otimes \text{Sym}^{k - q - 1}\mathcal{F})
\Rightarrow
H_{p + q}(\mathcal{C}, \text{Sym}^k(\mathcal{F}))
$$
By Cohomology on Sites, Lemma \ref{sites-cohomology-lemma-eilenberg-zilber}
we have
$$
L\pi_!(\wedge^{q + 1}\mathcal{F} \otimes \text{Sym}^{k - q - 1}\mathcal{F}) =
L\pi_!(\wedge^{q + 1}\mathcal{F}) \otimes_B^\mathbf{L}
L\pi_!(\text{Sym}^{k - q - 1}\mathcal{F}))
$$
It follows (from the construction of derived tensor products) that
the induction hypothesis combined with the vanishing
of $H_0(\mathcal{C}, \wedge^{q + 1}(\mathcal{F})) = 0$ will prove what we want.
This is true because $\wedge^{q + 1}(\mathcal{F})$ is a quotient
of $\mathcal{F}^{\otimes q + 1}$ and
$H_0(\mathcal{C}, \mathcal{F}^{\otimes q + 1})$
is a quotient of $H_0(\mathcal{C}, \mathcal{F})^{\otimes q + 1}$
which is zero.
\end{proof}

\begin{remark}
\label{remark-first-homology-symmetric-power}
In the situation of Lemma \ref{lemma-vanishing-symmetric-powers}
one can show that
$H_k(\mathcal{C}, \text{Sym}^k(\mathcal{F})) =
\wedge^k_B(H_1(\mathcal{C}, \mathcal{F}))$.
Namely, it can be deduced from the proof that
$H_k(\mathcal{C}, \text{Sym}^k(\mathcal{F}))$ is the $S_k$-coinvariants
of
$$
H^{-k}(L\pi_!(\mathcal{F}) \otimes_B^\mathbf{L}
L\pi_!(\mathcal{F}) \otimes_B^\mathbf{L}
\ldots \otimes_B^\mathbf{L} L\pi_!(\mathcal{F})) =
H_1(\mathcal{C}, \mathcal{F})^{\otimes k}
$$
Thus our claim is that this action is given by the usual action
of $S_k$ on the tensor product multiplied by the sign character.
To prove this one has to work through the sign conventions
in the definition of the total complex associated to a
multi-complex. We omit the verification.
\end{remark}

\begin{lemma}
\label{lemma-map-tors-zero}
Let $A$ be a ring. Let $P = A[E]$ be a polynomial ring.
Set $I = (e; e \in E) \subset P$. The maps
$\text{Tor}_i^P(A, I^{n + 1}) \to \text{Tor}_i^P(A, I^n)$
are zero for all $i$ and $n$.
\end{lemma}

\begin{proof}
Denote $x_e \in P$ the variable corresponding to $e \in E$.
A free resolution of $A$ over $P$ is given by the Koszul complex
$K_\bullet$ on the $x_e$. Here $K_i$ has basis given by wedges
$e_1 \wedge \ldots \wedge e_i$, $e_1, \ldots, e_i \in E$ and $d(e) = x_e$.
Thus $K_\bullet \otimes_P I^n = I^nK_\bullet$ computes
$\text{Tor}_i^P(A, I^n)$. Observe that everything is graded
with $\deg(x_e) = 1$, $\deg(e) = 1$, and $\deg(a) = 0$ for $a \in A$.
Suppose $\xi \in I^{n + 1}K_i$ is a cocycle homogeneous of degree $m$.
Note that $m \geq i + 1 + n$. Then $\xi = \text{d}\eta$ for some
$\eta \in K_{i + 1}$ as $K_\bullet$ is exact in degrees $ > 0$.
(The case $i = 0$ is left to the reader.)
Now $\deg(\eta) = m \geq i + 1 + n$. Hence writing $\eta$
in terms of the basis we see the coordinates are in $I^n$.
Thus $\xi$ maps to zero in the homology of $I^nK_\bullet$ as desired.
\end{proof}

\begin{proof}
Let $I \subset A$ be the kernel of $A \to B$. Let
$\mathcal{J} \subset \mathcal{O}$
be the kernel of $\mathcal{O} \to \underline{B}$. Then
$I\mathcal{O} \subset \mathcal{J}$. Set
$\mathcal{K} = \mathcal{J}/I\mathcal{O}$ and
$\overline{\mathcal{O}} = \mathcal{O}/I\mathcal{O}$.

\medskip\noindent
For every object $U = (P \to B)$ of $\mathcal{C}_{B/A}$
we can choose an isomorphism $P = A[E]$ such that the map
$P \to B$ maps each $e \in E$ to zero. Then
$J = \mathcal{J}(U) \subset P = \mathcal{O}(U)$
is equal to $J = IP + (e; e \in E)$. Moreover
$\overline{\mathcal{O}}(U) = B[E]$ and $K = \mathcal{K}(U) = (e; e \in E)$
is the ideal generated by the variables in the polynomial ring $B[E]$.
In particular it is clear that
$$
K/K^2 \xrightarrow{\text{d}} \Omega_{P/A} \otimes_P B
$$
is a bijection. In other words, $\Omega = \mathcal{K}/\mathcal{K}^2$
and $\text{Sym}_B^k(\Omega) = \mathcal{K}^k/\mathcal{K}^{k + 1}$.
Note that $\pi_!(\Omega) = \Omega_{B/A} = 0$ (Lemma \ref{lemma-identify-H0})
as $A \to B$ is surjective
(Algebra, Lemma \ref{algebra-lemma-trivial-differential-surjective}).
By Lemma \ref{lemma-vanishing-symmetric-powers} we conclude that
$$
H_i(\mathcal{C}_{B/A}, \mathcal{K}^k/\mathcal{K}^{k + 1}) =
H_i(\mathcal{C}_{B/A}, \text{Sym}^k_{\underline{B}}(\Omega)) = 0
$$
for $i < k$. This proves the final statement of the theorem.

\medskip\noindent
The approach to the theorem is to note that
$$
B \otimes_A^\mathbf{L} B = L\pi_!(\mathcal{O}) \otimes_A^\mathbf{L} B =
L\pi_!(\mathcal{O} \otimes_{\underline{A}}^\mathbf{L} \underline{B}) =
L\pi_!(\overline{\mathcal{O}})
$$
The first equality by Lemma \ref{lemma-apply-O-B-comparison},
the second equality by
Cohomology on Sites, Lemma \ref{sites-cohomology-lemma-change-of-rings}, and
the third equality as $\mathcal{O}$ is flat over $\underline{A}$.
The sheaf $\overline{\mathcal{O}}$ has a filtration
$$
\ldots \subset
\mathcal{K}^3 \subset
\mathcal{K}^2 \subset
\mathcal{K} \subset
\overline{\mathcal{O}}
$$
This induces a filtration $F$ on a complex $C$ representing
$L\pi_!(\overline{\mathcal{O}})$ with $F^pC$ representing
$L\pi_!(\mathcal{K}^p)$ (construction of $C$ and $F$ omitted).
Consider the spectral sequence of
Homology, Section \ref{homology-section-filtered-complex}
associated to $(C, F)$. It has $E_1$-page
$$
E_1^{p, q} = H_{- p - q}(\mathcal{C}_{B/A}, \mathcal{K}^p/\mathcal{K}^{p + 1})
\quad\Rightarrow\quad
H_{- p - q}(\mathcal{C}_{B/A}, \overline{\mathcal{O}}) = 
\text{Tor}_{- p - q}^A(B, B)
$$
and differentials $E_r^{p, q} \to E_r^{p + r, q - r + 1}$. To show convergence
we will show that for every $k$ there exists a $c$ such that
$H_i(\mathcal{C}_{B/A}, \mathcal{K}^n) = 0$
for $i < k$ and $n > c$\footnote{A posteriori
the ``correct'' vanishing $H_i(\mathcal{C}_{B/A}, \mathcal{K}^n) = 0$ for
$i < n$ can be concluded.}.

\medskip\noindent
Given $k \geq 0$ set $c = k^2$. We claim that
$$
H_i(\mathcal{C}_{B/A}, \mathcal{K}^{n + c}) \to
H_i(\mathcal{C}_{B/A}, \mathcal{K}^n)
$$
is zero for $i < k$ and all $n \geq 0$. Note that
$\mathcal{K}^n/\mathcal{K}^{n + c}$ has a finite filtration whose successive
quotients $\mathcal{K}^m/\mathcal{K}^{m + 1}$, $n \leq m < n + c$
have $H_i(\mathcal{C}_{B/A}, \mathcal{K}^m/\mathcal{K}^{m + 1}) = 0$
for $i < n$ (see above). Hence the claim implies
$H_i(\mathcal{C}_{B/A}, \mathcal{K}^{n + c}) = 0$ for $i < k$ and all
$n \geq k$ which is what we need to show.

\medskip\noindent
Proof of the claim. Recall that for any $\mathcal{O}$-module $\mathcal{F}$
the map $\mathcal{F} \to \mathcal{F} \otimes_\mathcal{O}^\mathbf{L} B$
induces an isomorphism on applying $L\pi_!$, see
Lemma \ref{lemma-O-homology-B-homology}.
Consider the map
$$
\mathcal{K}^{n + k} \otimes_\mathcal{O}^\mathbf{L} B
\longrightarrow
\mathcal{K}^n \otimes_\mathcal{O}^\mathbf{L} B
$$
We claim that this map induces the zero map on cohomology sheaves
in degrees $0, -1, \ldots, - k + 1$. If this second claim holds, then
the $k$-fold composition
$$
\mathcal{K}^{n + c} \otimes_\mathcal{O}^\mathbf{L} B
\longrightarrow
\mathcal{K}^n \otimes_\mathcal{O}^\mathbf{L} B
$$
factors through $\tau_{\leq -k}\mathcal{K}^n \otimes_\mathcal{O}^\mathbf{L} B$
hence induces zero on $H_i(\mathcal{C}_{B/A}, -) = L_i\pi_!( - )$
for $i < k$, see
Derived Categories, Lemma \ref{derived-lemma-trick-vanishing-composition}.
By the remark above this means the same thing is true for
$H_i(\mathcal{C}_{B/A}, \mathcal{K}^{n + c}) \to
H_i(\mathcal{C}_{B/A}, \mathcal{K}^n)$
which proves the (first) claim.

\medskip\noindent
Proof of the second claim. The statement is local, hence we may work
over an object $U = (P \to B)$ as above. We have to show
the maps
$$
\text{Tor}_i^P(B, K^{n + k}) \to \text{Tor}_i^P(B, K^n)
$$
are zero for $i < k$. There is a spectral sequence
$$
\text{Tor}_a^P(P/IP, \text{Tor}_b^{P/IP}(B, K^n))
\Rightarrow
\text{Tor}_{a + b}^P(B, K^n),
$$
see More on Algebra, Example \ref{more-algebra-example-tor-change-rings}.
Thus it suffices to prove the maps
$$
\text{Tor}_i^{P/IP}(B, K^{n + 1}) \to \text{Tor}_i^{P/IP}(B, K^n)
$$
are zero for all $i$. This is Lemma \ref{lemma-map-tors-zero}.
\end{proof}

\bibliographystyle{plain}
\bibliography{my}
\end{document}
